\chapter{Sieves}

\begin{center}\itshape
The sieve of Eratosthenes is the oldest algorithm in number theory, and for two thousand years it was also the only one that counted primes.
\end{center}

Around 240 BC, Eratosthenes of Cyrene, who also measured the circumference of the Earth, described how to find the primes: write the numbers, strike out the
multiples of $2$, then of $3$, then of $5$, and so on; what survives the multiples of every prime up to $\sqrt x$ is the primes up to $x$. It is an algorithm, not a theorem,
and it gives no formula. In 1808 Legendre turned it into one. The number of survivors is an inclusion--exclusion sum, exact, and it looks like it should give the prime
number theorem at once. It does not: the formula has $2^{\pi(\sqrt x)}$ terms, and the errors pile up beyond all control.

For another century nobody could make a sieve count anything. Then, between 1915 and 1919, a young Norwegian, Viggo Brun, found how to stop the inclusion--exclusion
early, with an error that could be bounded. He applied it to the oldest unsolved question about primes: are there infinitely many twin primes, like $11,13$ or $101,103$? He
could not answer it, and nobody can yet. But he proved something no one had suspected: the twin primes are so rare that the sum of their reciprocals converges,
\[
\begin{aligned}\Bigl(\frac13+\frac15\Bigr)&+\Bigl(\frac15+\frac17\Bigr)\eqbr+\Bigl(\frac1{11}+\frac1{13}\Bigr)+\cdots=B_2<\infty,\end{aligned}
\]
whereas $\sum\frac1p$ diverges (Chapter~1). The value $B_2=1.9021605\ldots$ has its own place in history: in 1994 Thomas Nicely, computing it on a new Pentium processor,
noticed that the chip divided incorrectly. Intel replaced the processors, at a cost of about $475$ million dollars.

In 1947 Atle Selberg replaced Brun's combinatorics by a single idea: an upper bound for a count is a sum of squares, and the best sum of squares is the solution of a quadratic
optimisation problem. Six years earlier, Yuri Linnik had found a sieve of a different kind, the \emph{large sieve}, which removes many residue classes for each prime and is really
a statement about harmonic analysis. This chapter proves both, with their first applications: twin primes, primes in short intervals, and sets that look like squares.
Selberg's weights are still at the heart of the subject: they are the starting point of Zhang's and Maynard's proofs (2013) that there are infinitely many pairs of primes at
bounded distance.

\begin{think}
If divisibility by different primes were independent events, the proportion of $n\le x$ with no prime factor $\le z$ would be $\prod_{p\le z}(1-\frac1p)$. For $z=\sqrt x$, the survivors
are exactly $1$ and the primes in $]\sqrt x,x]$. By Mertens's theorem and the prime number theorem, compare the two. Which is wrong, and by how much?
\end{think}

\section{Legendre's sieve, and why it fails}

Let $P(z)=\prod_{p\le z}p$ and let $\Phi(x,z)$ be the number of integers $1\le n\le x$ with no prime factor $\le z$. Since $\sum_{d\mid m}\mu(d)$ is $1$ if $m=1$ and $0$ otherwise (Volume~2),
\[
\Phi(x,z)=\sum_{n\le x}\sum_{d\mid(n,P(z))}\mu(d)=\sum_{d\mid P(z)}\mu(d)\Bigl\lfloor\frac xd\Bigr\rfloor.
\]
This is Legendre's identity: inclusion--exclusion over the primes $\le z$. Writing $\lfloor\frac xd\rfloor=\frac xd-\{\frac xd\}$ separates a main term from an error.

\begin{proposition}[Legendre, 1808]\label{prop:legendre}
For $x\ge1$ and $z\ge2$,
\[
\Phi(x,z)=x\prod_{p\le z}\Bigl(1-\frac1p\Bigr)+E,\qquad|E|\le2^{\pi(z)}.
\]
\end{proposition}
\begin{proof}
$\sum_{d\mid P(z)}\mu(d)\frac xd=x\prod_{p\le z}(1-\frac1p)$ by expanding the product, and each of the $2^{\pi(z)}$ fractional parts is in $[0,1[$.
\end{proof}

The main term is right only if the error is small, so $z$ must be tiny: $2^{\pi(z)}\le\sqrt x$ forces $z$ to be about $\log x$. With $z=\log x$, Mertens's theorem
(Theorem~\ref{thm:mertens3}) gives $\pi(x)\le z+\Phi(x,z)=O(\frac x{\log\log x})$ (Exercise~\ref{ex:legendrepi}): true, and very weak.

Worse, the main term is \emph{wrong} at $z=\sqrt x$. There $\Phi(x,\sqrt x)=\pi(x)-\pi(\sqrt x)+1\sim\frac x{\log x}$, while $x\prod_{p\le\sqrt x}(1-\frac1p)\sim\frac{2e^{-\gamma}x}{\log x}=\frac{1.1229\,x}{\log x}$. The
$2^{\pi(\sqrt x)}$ fractional parts, each in $[0,1[$, conspire to remove $12\%$ of the main term (Exercise~\ref{ex:legendrefail}). This is the factor already met in Chapter~1, and it is not a
defect of Legendre's bookkeeping: a sieve that only knows how many $n$ lie in each residue class modulo small numbers cannot tell the primes from integers with exactly two large prime
factors. Selberg made this precise as the \emph{parity problem}: sieve methods alone cannot distinguish integers with an even number of prime factors from those with an odd number.
That is why sieves give upper bounds of the right order of magnitude, and lower bounds only for ``almost primes''.

\section{Selberg's sieve}

\subsection*{The setting}

Let $\mathcal A$ be a finite list of positive integers (repetitions allowed), $\mathcal P$ a set of primes, $z\ge2$, and $P(z)=\prod_{p\in\mathcal P,\,p\le z}p$. We want an upper bound for
\[
S(\mathcal A,\mathcal P,z)=\#\{a\in\mathcal A:\gcd(a,P(z))=1\}.
\]
For $d\mid P(z)$ let $\mathcal A_d$ be the sublist of multiples of $d$. The only information we use is that $|\mathcal A_d|$ is roughly a proportion $g(d)$ of a total $X$:
\[
|\mathcal A_d|=g(d)X+r_d\qquad(d\mid P(z)),
\]
where $g$ is multiplicative, $0<g(p)<1$ for $p\in\mathcal P$, and $r_d$ is a remainder we hope is small. For the integers in an interval of length $X$, $g(d)=\frac1d$ and $|r_d|\le1$. Let
\[
h(p)=\frac{g(p)}{1-g(p)},\qquad G(y)=\sum_{\substack{d\le y\\d\mid P(z)}}\mu^2(d)h(d),
\]
with $h$ extended multiplicatively to squarefree $d$.

\begin{theorem}[Selberg, 1947]\label{thm:selberg}
For every $y\ge1$,
\[
S(\mathcal A,\mathcal P,z)\le\frac X{G(y)}+\sum_{\substack{d_1,d_2\le y\\d_1,d_2\mid P(z)}}|r_{[d_1,d_2]}|,
\]
where $[d_1,d_2]$ is the least common multiple.
\end{theorem}

The idea fits in one line. If $\lambda_1=1$, then for every integer $m$, $\bigl(\sum_{d\mid m}\lambda_d\bigr)^2$ is $1$ when $m=1$ and \emph{at least $0$} otherwise. So any real weights with $\lambda_1=1$ give an upper
bound, and we choose the weights that make the bound smallest. Brun had to engineer signs; Selberg's squares are positive for free.

\begin{proof}
\emph{Squares.} Let $(\lambda_d)$ be real numbers indexed by the $d\mid P(z)$ with $d\le y$, with $\lambda_1=1$. For $a\in\mathcal A$, $\bigl(\sum_{d\mid(a,P(z))}\lambda_d\bigr)^2$ is $1$ if $\gcd(a,P(z))=1$, and is $\ge0$ otherwise. So
\[
\begin{aligned}S&\le\sum_{a\in\mathcal A}\Bigl(\sum_{d\mid(a,P(z))}\lambda_d\Bigr)^2=\sum_{d_1,d_2}\lambda_{d_1}\lambda_{d_2}|\mathcal A_{[d_1,d_2]}|\\&=XQ+\sum_{d_1,d_2}\lambda_{d_1}\lambda_{d_2}r_{[d_1,d_2]},\end{aligned}
\]
since $d_1\mid a$ and $d_2\mid a$ means $[d_1,d_2]\mid a$, and where $Q=\sum_{d_1,d_2}\lambda_{d_1}\lambda_{d_2}g([d_1,d_2])$. We minimise the quadratic form $Q$, then check that the minimiser has $|\lambda_d|\le1$, which bounds the
second sum by the remainder of the theorem.

\emph{Diagonalising $Q$.} For squarefree $d_1,d_2$, $g([d_1,d_2])=\frac{g(d_1)g(d_2)}{g((d_1,d_2))}$, by multiplicativity. Since $\frac1{g(p)}=1+\frac1{h(p)}$, for squarefree $m$ we have $\frac1{g(m)}=\prod_{p\mid m}\bigl(1+\frac1{h(p)}\bigr)=\sum_{e\mid m}\frac1{h(e)}$.
Hence
\[
\begin{aligned}Q&=\sum_{d_1,d_2}\lambda_{d_1}g(d_1)\lambda_{d_2}g(d_2)\eqbr\quad\times\sum_{e\mid(d_1,d_2)}\frac1{h(e)}=\sum_e\frac{u_e^2}{h(e)},\end{aligned}
\]
where $u_e=\sum_{e\mid d}\lambda_dg(d)$, with $e$ and $d$ running over the divisors of $P(z)$ up to $y$. This change of variables is invertible: $\lambda_dg(d)=\sum_{d\mid e}\mu(\frac ed)u_e$, as one checks by substituting the definition of $u_e$ and using that $\sum_{d\mid e\mid k}\mu(\frac ed)$ is $1$ if $k=d$ and $0$ otherwise. The
constraint $\lambda_1=1$ becomes $\sum_e\mu(e)u_e=1$.

\emph{The minimum.} By Cauchy--Schwarz,
\[
\begin{aligned}1=\Bigl(\sum_e\mu(e)u_e\Bigr)^2&\le\Bigl(\sum_e\frac{u_e^2}{h(e)}\Bigr)\Bigl(\sum_eh(e)\Bigr)\eqbr=Q\cdot G(y),\end{aligned}
\]
with equality when $u_e=\frac{\mu(e)h(e)}{G(y)}$. With this choice, $Q=\frac1{G(y)}$.

\emph{The weights are at most $1$.} Write $e=dm$ in the inversion formula; $m$ is squarefree and coprime to $d$, $\mu(m)\mu(dm)=\mu(d)\mu(m)^2=\mu(d)$, and $h(dm)=h(d)h(m)$. So
\[
\lambda_d=\frac{\mu(d)h(d)}{g(d)G(y)}\sum_{\substack{m\le y/d,\ m\mid P(z)\\(m,d)=1}}\mu^2(m)h(m).
\]
Now $\frac{h(d)}{g(d)}=\prod_{p\mid d}\frac1{1-g(p)}=\prod_{p\mid d}(1+h(p))=\sum_{k\mid d}h(k)$. Each product $km$ with $k\mid d$ and $m$ in the sum is a squarefree divisor of $P(z)$, at most $y$, and different pairs give
different products ($k=(km,d)$). So $\bigl(\sum_{k\mid d}h(k)\bigr)\bigl(\sum_m\mu^2(m)h(m)\bigr)\le G(y)$, which is $|\lambda_d|\le1$. Then $|\lambda_{d_1}\lambda_{d_2}r_{[d_1,d_2]}|\le|r_{[d_1,d_2]}|$.
\end{proof}

To use the theorem we need a lower bound for $G(y)$. For the integers themselves it is a pleasant identity in disguise.

\begin{lemma}\label{lem:Glog}
For $y\ge1$, $\displaystyle\sum_{d\le y}\frac{\mu^2(d)}{\varphi(d)}\ge\log y$.
\end{lemma}
\begin{proof}
For squarefree $d$, $\frac1{\varphi(d)}=\frac1d\prod_{p\mid d}\bigl(1-\frac1p\bigr)^{-1}=\frac1d\prod_{p\mid d}\sum_{k\ge0}p^{-k}=\sum_{\operatorname{rad}(m)=d}\frac1m$, where $\operatorname{rad}(m)$ is the product of the primes dividing $m$. Every $m\le y$
has $\operatorname{rad}(m)\le m\le y$, so the sum is at least $\sum_{m\le y}\frac1m\ge\log y$.
\end{proof}

\subsection*{Primes in short intervals}

The first application counts primes in any interval of length $N$, wherever it is. The prime number theorem says nothing about $[x,x+N]$ when $N$ is much smaller than $x$; the sieve
does not care where the interval is.

\begin{theorem}[Brun--Titchmarsh]\label{thm:bt}
Uniformly for $x\ge0$ and $N\ge3$,
\[
\pi(x+N)-\pi(x)\le\frac{2N}{\log N}\Bigl(1+O\Bigl(\frac{\log\log N}{\log N}\Bigr)\Bigr).
\]
\end{theorem}
\begin{proof}
Let $\mathcal A$ be the integers in $]x,x+N]$, $\mathcal P$ all primes, and $X=N$. For $d\ge1$, $|\mathcal A_d|=\frac Nd+r_d$ with $|r_d|\le1$; so $g(d)=\frac1d$ and $h(p)=\frac1{p-1}$, $h(d)=\frac1{\varphi(d)}$. Take
$y=z=\frac{\sqrt N}{\log N}$ (for $N$ large; small $N$ go into the $O$). Every squarefree $d\le y$ divides $P(z)$, so by Lemma~\ref{lem:Glog}, $G(y)\ge\log y=\frac12\log N-\log\log N$. The remainder
is at most the number of pairs $(d_1,d_2)$, that is $y^2=\frac N{\log^2N}$. A prime in $]x,x+N]$ is either $\le z$ or coprime to $P(z)$, so
\[
\begin{aligned}\pi(x+N)-\pi(x)&\le z+\frac N{\frac12\log N-\log\log N}\eqbr+\frac N{\log^2N},\end{aligned}
\]
which is the claim.
\end{proof}

By the prime number theorem, an interval $]x,x+N]$ with $N$ about $x$ contains about $\frac N{\log x}$ primes. Brun--Titchmarsh bounds every interval of length $N$ by $\frac{2N}{\log N}$, wherever the interval sits. If $x=N^k$, the ratio of this bound to $\frac N{\log x}$ is $2k$, not $2$. Montgomery and Vaughan (1973) removed the $O$ term: $\pi(x+N)-\pi(x)\le\frac{2N}{\log N}$ for all $N\ge2$. The constant $2$ that would bear on Siegel zeros is the constant in the arithmetic-progression inequality $\pi(x;q,a)\le\frac{2x}{\varphi(q)\log(x/q)}$. Chapter~3 names Siegel zeros and does not prove that inequality. The parity problem blocks sieve upper bounds of shape $2-\varepsilon$ both for short intervals and for arithmetic progressions. Improving the short-interval constant from $2$ to $2-\varepsilon$ would not rule out Siegel zeros.

\subsection*{Twin primes}

Let $\pi_2(x)$ be the number of primes $p\le x$ such that $p+2$ is also prime. To catch both, sieve the products $n(n+2)$: both factors are prime and larger than $z$ exactly when $n(n+2)$
has no prime factor $\le z$ (and $n>z$).

\begin{theorem}[Brun, 1919]\label{thm:twins}
$\pi_2(x)=O\Bigl(\dfrac x{\log^2x}\Bigr)$.
\end{theorem}
\begin{proof}
\emph{The data.} Let $\mathcal A=\{n(n+2):1\le n\le x\}$ and $\mathcal P$ all primes. For squarefree $d$, $d\mid n(n+2)$ is a condition on $n\bmod d$; by the Chinese remainder theorem, the number $\rho(d)$ of solutions modulo
$d$ is multiplicative, with $\rho(2)=1$ and $\rho(p)=2$ for $p$ odd (the classes $0$ and $-2$). Counting $n\le x$ in $\rho(d)$ classes modulo $d$, $|\mathcal A_d|=\frac{\rho(d)}dx+r_d$ with $|r_d|\le\rho(d)$. So
$g(d)=\frac{\rho(d)}d$, $h(2)=1$, and $h(p)=\frac{2}{p-2}$ for odd $p$.

\emph{The main term.} For odd $p$, $h(p)=\frac2p\cdot\frac1{1-2/p}=\sum_{k\ge1}\frac{2^k}{p^k}$. So for odd squarefree $d$, $h(d)=\sum_{\operatorname{rad}(m)=d}\frac{2^{\Omega(m)}}m\ge\sum_{\operatorname{rad}(m)=d}\frac{\tau(m)}m$, where
$\Omega(m)$ counts prime factors with multiplicity and $\tau(m)\le2^{\Omega(m)}$. As in Lemma~\ref{lem:Glog}, taking $z\ge y$,
\[
\begin{aligned}G(y)&\ge\sum_{\substack{m\le y\\m\text{ odd}}}\frac{\tau(m)}m=\sum_{\substack{ab\le y\\a,b\text{ odd}}}\frac1{ab}\ge\Bigl(\sum_{\substack{a\le\sqrt y\\a\text{ odd}}}\frac1a\Bigr)^2.\end{aligned}
\]
The odd $a\le Y$ dominate term by term the even $b\le Y+1$, so the last sum is at least $\frac12\log\frac{\sqrt y}2$, and $G(y)\ge\frac1{16}\log^2y\,(1+o(1))$.

\emph{The remainder.} $\rho([d_1,d_2])\le2^{\omega([d_1,d_2])}\le2^{\omega(d_1)}2^{\omega(d_2)}\le\tau(d_1)\tau(d_2)$, so the remainder is at most $\bigl(\sum_{d\le y}\tau(d)\bigr)^2$. And $\sum_{d\le y}\tau(d)=\sum_{a\le y}\lfloor\frac ya\rfloor\le y(1+\log y)$.

\emph{The choice.} Take $y=z=\frac{\sqrt x}{\log^3x}$. Then the remainder is at most $\frac x{\log^4x}$ for large $x$, $\log y\sim\frac12\log x$, and
\[
\begin{aligned}\pi_2(x)&\le z+\frac{16x}{\log^2y}(1+o(1))+\frac x{\log^4x}\eqbr=O\Bigl(\frac x{\log^2x}\Bigr).\end{aligned}
\]
\end{proof}

\begin{corollary}[Brun, 1919]\label{cor:brun}
The series $\sum\frac1p$ over the primes $p$ such that $p+2$ is prime converges.
\end{corollary}
\begin{proof}
By Lemma~\ref{lem:abelint} with $a_p=1$ for these $p$ and $f(t)=\frac1t$, $\sum_{p\le x}\frac1p=\frac{\pi_2(x)}x+\int_2^x\frac{\pi_2(t)}{t^2}\dd t$, and $\int_2^\infty\frac{\mathrm dt}{t\log^2t}<\infty$.
\end{proof}

\begin{figure}[ht]
\centering
\fig{ch08-twins}
\caption{The count $\pi_2(x)$ of twin prime pairs (charcoal) for $x\le10^5$, against the Hardy--Littlewood prediction $2C_2\int_2^x\frac{\mathrm dt}{\log^2t}$ (coral, dashed), with $C_2=0.6601618\ldots$ At $10^7$ the
prediction is $58\,754$, and $\pi_2(10^7)=58\,980$.}
\label{fig:twins}
\end{figure}

\begin{remark}[how many twins?]
The constant in our proof is about $64$; a more careful evaluation of $G(y)$ gives Selberg's $\pi_2(x)\le(16C_2+o(1))\frac x{\log^2x}$, where $C_2=\prod_{p>2}\bigl(1-\frac1{(p-1)^2}\bigr)=0.6601618\ldots$ is the \emph{twin prime constant}.
Hardy and Littlewood conjectured in 1923 that $\pi_2(x)\sim2C_2\frac x{\log^2x}$, so this bound is $8$ times too large. Sifting the shifted primes $p+2$ instead of the integers $n(n+2)$, with the Bombieri--Vinogradov theorem (Remark~\ref{rem:bv} below) to control the remainders, Bombieri and Davenport (1966) brought the factor down to $4$; the best known factor is about $3.4$ (Wu, 2004). The conjecture comes from the same
probabilistic heuristic as the think box, corrected prime by prime: Figure~\ref{fig:twins} and Exercise~\ref{ex:twinheuristic} explain it, and Chapter~10 returns to it. As for a lower bound,
the parity problem blocks the sieve; the best result is Chen's theorem (1973): there are infinitely many primes $p$ such that $p+2$ is a prime or a product of two primes.
\end{remark}

\begin{beauty}
Zhang (2013) proved that there are infinitely many pairs of primes at distance at most $70\,000\,000$. Months later Maynard, and independently Tao, found a simpler and stronger method:
for every $m$, there are bounded intervals containing $m$ primes infinitely often. Their tool is a Selberg sieve in several variables, with weights $\lambda_{d_1,\dots,d_k}$ chosen,
exactly as above, to optimise a quadratic form. A collaborative project (Polymath, 2014) brought the gap down to $246$. The twin prime conjecture is the gap $2$.
\end{beauty}

\section{The large sieve}

Selberg's sieve removes a few residue classes, one or two per prime. Sometimes we know much more: a set of integers that avoids \emph{half} of the classes modulo every prime, like
the squares. Linnik discovered in 1941 that such a set must be tiny. The modern proof, due to Rényi, Roth, Bombieri, Davenport, Halberstam, Gallagher and Montgomery, is pure harmonic analysis:
a set that is lopsided modulo $p$ has a large Fourier coefficient at the fractions $\frac ap$, and a trigonometric polynomial cannot be large at too many well-spaced points.

Write $e(\alpha)=e^{2\pi i\alpha}$ and $\|\alpha\|$ for the distance from $\alpha$ to the nearest integer. Points $\alpha_1,\dots,\alpha_R$ are \emph{$\delta$-spaced} if $\|\alpha_r-\alpha_s\|\ge\delta$ for $r\ne s$.

\begin{lemma}[Gallagher]\label{lem:gallagher}
Let $\delta>0$ and $f$ of class $C^1$ on $I=[\alpha-\frac\delta2,\alpha+\frac\delta2]$, with complex values. Then
\[
|f(\alpha)|\le\frac1\delta\int_I|f|+\frac12\int_I|f'|.
\]
\end{lemma}
\begin{proof}
For $\beta\in I$, $f(\alpha)=f(\beta)-\int_\alpha^\beta f'$. Average over $\beta\in I$: $|f(\alpha)|\le\frac1\delta\int_I|f|+\frac1\delta\int_I\int_{[\alpha,\beta]}|f'(t)|\dd t\dd\beta$. In the double integral, a point $t>\alpha$ is counted for the $\beta\in[t,\alpha+\frac\delta2]$,
a set of length at most $\frac\delta2$, and similarly for $t<\alpha$. So the double integral is at most $\frac\delta2\int_I|f'|$.
\end{proof}

\begin{theorem}[the large sieve inequality]\label{thm:largesieve}
Let $a_n$ be complex numbers for the integers $M<n\le M+N$, $S(\alpha)=\sum_na_ne(n\alpha)$, and let $\alpha_1,\dots,\alpha_R$ be $\delta$-spaced with $0<\delta\le\frac12$. Then
\[
\sum_{r=1}^R|S(\alpha_r)|^2\le\Bigl(\pi N+\frac1\delta\Bigr)\sum_n|a_n|^2.
\]
\end{theorem}
\begin{proof}
Let $c=M+\frac{N+1}2$, the centre of the range, and $T(\alpha)=e(-c\alpha)S(\alpha)=\sum_na_ne((n-c)\alpha)$, so $|T|=|S|$ and $|n-c|\le\frac{N-1}2$. Apply Lemma~\ref{lem:gallagher} to $f=T^2$, with $f'=2TT'$, at each $\alpha_r$.
The intervals of length $\delta$ centred at the $\alpha_r$ do not overlap modulo $1$, and $|T|$ and $|T'|$ are $1$-periodic, so
\[
\sum_r|S(\alpha_r)|^2\le\frac1\delta\int_0^1|T|^2+\int_0^1|T||T'|.
\]
By orthogonality, $\int_0^1e((n-c)\alpha)\overline{e((m-c)\alpha)}\dd\alpha=\int_0^1e((n-m)\alpha)\dd\alpha$ is $1$ if $n=m$ and $0$ otherwise. So $\int_0^1|T|^2=\sum|a_n|^2$, and since $T'(\alpha)=\sum2\pi i(n-c)a_ne((n-c)\alpha)$,
$\int_0^1|T'|^2=4\pi^2\sum(n-c)^2|a_n|^2\le\pi^2N^2\sum|a_n|^2$. Cauchy--Schwarz bounds $\int|T||T'|$ by $\pi N\sum|a_n|^2$.
\end{proof}

The constant is not optimal: Montgomery and Vaughan (1973), and Selberg, obtained $N-1+\frac1\delta$, which is sharp. What matters is the shape: $N$ is the number of coefficients, $\frac1\delta$ the
number of points one can fit, and the inequality says these two effects add rather than multiply.

The fractions $\frac aq$ with $q\le Q$ and $\gcd(a,q)=1$, $1\le a\le q$, are distinct modulo $1$ and $\frac1{Q^2}$-spaced: $\bigl|\frac aq-\frac{a'}{q'}\bigr|=\frac{|aq'-a'q|}{qq'}\ge\frac1{qq'}$ when they differ. So, for $Q\ge2$,
\begin{equation}\label{eq:arithls}
\sum_{q\le Q}\ \sum_{\substack{a=1\\(a,q)=1}}^q\Bigl|S\Bigl(\frac aq\Bigr)\Bigr|^2\le(\pi N+Q^2)\sum_n|a_n|^2.
\end{equation}

Now the sieve. A set that misses residue classes modulo $p$ is unevenly spread modulo $p$, and unevenness shows up in the Fourier coefficients.

\begin{lemma}\label{lem:lsprime}
Let $\mathcal A$ be a finite set of integers, $c_n$ complex numbers for $n\in\mathcal A$, and $S(\alpha)=\sum_{n\in\mathcal A}c_ne(n\alpha)$. Suppose that for each prime $p$ in a set $\mathcal Q$, $\mathcal A$ meets at most $p-\omega(p)$ classes
modulo $p$, with $0\le\omega(p)<p$. Then for every squarefree $q$ whose prime factors lie in $\mathcal Q$,
\[
\sum_{\substack{a=1\\(a,q)=1}}^q\Bigl|S\Bigl(\frac aq\Bigr)\Bigr|^2\ge|S(0)|^2\prod_{p\mid q}\frac{\omega(p)}{p-\omega(p)}.
\]
\end{lemma}
\begin{proof}
\emph{A prime $q=p$.} Let $Z_h=\sum_{n\in\mathcal A,\,n\equiv h\,(p)}c_n$, so $S(\frac ap)=\sum_{h=0}^{p-1}Z_he(\frac{ah}p)$. By orthogonality of the characters $h\mapsto e(\frac{ah}p)$, $\sum_{a=0}^{p-1}|S(\frac ap)|^2=p\sum_h|Z_h|^2$. At most $p-\omega(p)$ of the
$Z_h$ are non-zero, so by Cauchy--Schwarz $|S(0)|^2=|\sum_hZ_h|^2\le(p-\omega(p))\sum_h|Z_h|^2$. Removing the term $a=0$:
\[
\begin{aligned}\sum_{a=1}^{p-1}\Bigl|S\Bigl(\frac ap\Bigr)\Bigr|^2&\ge\Bigl(\frac p{p-\omega(p)}-1\Bigr)|S(0)|^2\eqbr=\frac{\omega(p)}{p-\omega(p)}|S(0)|^2.\end{aligned}
\]
\emph{Induction.} Let $q=q_1p$ with $p\nmid q_1$. By the Chinese remainder theorem, as $a_1$ runs over the reduced classes modulo $q_1$ and $a_2$ over those modulo $p$, $\frac{a_1}{q_1}+\frac{a_2}p$ runs over the
reduced fractions $\frac aq$ modulo $1$. For fixed $a_1$, apply the prime case to the coefficients $c_ne(\frac{na_1}{q_1})$, which live on the same set $\mathcal A$: $\sum_{a_2}|S(\frac{a_1}{q_1}+\frac{a_2}p)|^2\ge\frac{\omega(p)}{p-\omega(p)}|S(\frac{a_1}{q_1})|^2$.
Sum over $a_1$ and use the induction hypothesis for $q_1$.
\end{proof}

\begin{theorem}[the arithmetic large sieve, Montgomery 1968]\label{thm:arithls}
Let $\mathcal A$ be a set of integers in $]M,M+N]$, $Q\ge2$, and suppose that for every prime $p\le Q$, $\mathcal A$ meets at most $p-\omega(p)$ classes modulo $p$. Then
\[
\begin{gathered}|\mathcal A|\le\frac{\pi N+Q^2}L,\qq L=\sum_{q\le Q}\mu^2(q)\prod_{p\mid q}\frac{\omega(p)}{p-\omega(p)}.\end{gathered}
\]
\end{theorem}
\begin{proof}
Take $c_n=1$ on $\mathcal A$, so $S(0)=|\mathcal A|$. Sum Lemma~\ref{lem:lsprime} over squarefree $q\le Q$ (the term $q=1$ is $|S(0)|^2$ on both sides) and compare with \eqref{eq:arithls}: $|\mathcal A|^2L\le(\pi N+Q^2)|\mathcal A|$.
\end{proof}

\begin{proposition}\label{prop:squares}
Let $\mathcal A\subset[1,N]$ be a set of integers such that, for every prime $p\le\sqrt N$, every element of $\mathcal A$ is a square modulo $p$. Then $|\mathcal A|=O(\sqrt N\log N)$.
\end{proposition}
\begin{proof}
For odd $p$ there are $\frac{p+1}2$ squares modulo $p$, counting $0$ (Volume~3, Chapter~11), so $\mathcal A$ misses $\omega(p)=\frac{p-1}2$ classes, and $\frac{\omega(p)}{p-\omega(p)}=\frac{p-1}{p+1}\ge\frac12$. Keeping only the prime terms $q=p$ in $L$,
$L\ge\frac12(\pi(Q)-1)$. With $Q=\sqrt N$ and $\pi(Q)\sim\frac Q{\log Q}$ (Chapter~5), Theorem~\ref{thm:arithls} gives $|\mathcal A|\le\frac{(\pi+1)N}{L}=O(\sqrt N\log N)$.
\end{proof}

The squares themselves satisfy the hypothesis and number $\sqrt N$, so the bound is sharp up to the logarithm (Gallagher's ``larger sieve'' removes it). Is anything else in $\mathcal A$? The
one finds numerically that the integers $n\le10^6$ that are squares modulo every prime $p\le1000$ are exactly the $1000$ squares. Whether a large set that looks like the squares modulo
every prime must essentially be the values of a quadratic polynomial is an open problem, the \emph{inverse sieve problem}.

\begin{remark}[Bombieri--Vinogradov]\label{rem:bv}
The large sieve has a version with Dirichlet characters, obtained from Theorem~\ref{thm:largesieve} through Gauss sums (Chapter~3). Its most famous consequence is the Bombieri--Vinogradov theorem (1965):
let $E(x;q)=\max_{(a,q)=1}\bigl|\psi(x;q,a)-\frac x{\varphi(q)}\bigr|$; for every $A>0$ there is $B$ such that
\[
\sum_{q\le\sqrt x/\log^Bx}E(x;q)=O\Bigl(\frac x{\log^Ax}\Bigr).
\]
On average over the moduli, primes in progressions are as well distributed as the generalized Riemann hypothesis predicts. This ``GRH on average'' is the fuel of Chen's theorem and of Zhang's
and Maynard's bounded gaps. See Davenport, Chapters~27--28.
\end{remark}

\section{Exercises}

\begin{exercise}[$\star$]\label{ex:legendrepi}
Use Proposition~\ref{prop:legendre} with $z=\log x$ to show that $\pi(x)=O\bigl(\frac x{\log\log x}\bigr)$.
\end{exercise}

\begin{exercise}[$\star$]\label{ex:legendrefail}
Show that $\sum_{d\mid P(\sqrt x)}\mu(d)\bigl\{\frac xd\bigr\}\sim(2e^{-\gamma}-1)\frac x{\log x}$. So the fractional parts in Legendre's formula, which look like noise, add up to a positive bias of $12\%$ of the main term.
\end{exercise}

\begin{exercise}[$\star$]\label{ex:sophie}
A prime $p$ is a Sophie Germain prime if $2p+1$ is also prime. Show that there are $O\bigl(\frac x{\log^2x}\bigr)$ of them up to $x$.
\end{exercise}

\begin{exercise}[$\star\star$]\label{ex:btap}
Let $1\le q<x$ and $\gcd(a,q)=1$. Show that $\pi(x;q,a)\le\frac{2x}{\varphi(q)\log(x/q)}\bigl(1+O\bigl(\frac{\log\log(x/q)}{\log(x/q)}\bigr)\bigr)$, uniformly in $q$ and $a$.
\end{exercise}

\begin{exercise}[$\star\star$]\label{ex:goldbachup}
For even $N\ge4$, let $r(N)$ be the number of primes $p\le N$ such that $N-p$ is prime. Show that $r(N)=O\Bigl(\bigl(\frac N{\varphi(N)}\bigr)^2\frac N{\log^2N}\Bigr)$.
\end{exercise}

\begin{exercise}[$\star\star$, the signature exercise]\label{ex:twinheuristic}
(a) Show that the sum of $\frac1p+\frac1{p+2}$ over twin pairs with $p>x$ is $O(\frac1{\log x})$.
(b) Explain the Hardy--Littlewood prediction: for $n$ random near $x$, the ``probability'' that $n$ and $n+2$ are both prime is $\frac1{\log^2x}$ if the events were independent; show that correcting for each
prime $p$ (the chance that $p\nmid n(n+2)$ is $1-\frac{\rho(p)}p$ instead of $(1-\frac1p)^2$) multiplies this by $2C_2$.
(c) Assuming $\pi_2(x)\sim2C_2\frac x{\log^2x}$, show that the tail in (a) is $\sim\frac{4C_2}{\log x}$, and use the partial sum $1.7383570$ up to $10^7$ to estimate Brun's constant $B_2$.
\end{exercise}

\begin{exercise}[$\star$]\label{ex:lsprimes}
Apply Theorem~\ref{thm:arithls} to the primes in $]M,M+N]$ to show $\pi(M+N)-\pi(M)\le\frac{2\pi N}{\log N}(1+o(1))$. Where is the constant lost compared with Theorem~\ref{thm:bt}?
\end{exercise}

\begin{exercise}[$\star\star\star$]\label{ex:nsq1}
Show that the number of $n\le x$ such that $n^2+1$ is prime is $O\bigl(\frac x{\log x}\bigr)$. (Whether it tends to infinity is Landau's fourth problem, open since 1912.)
\end{exercise}

\begin{exercise}[$\star$]\label{pb:bonferroni}
\emph{Bonferroni's inequalities.} Let $A_1,\dots,A_m$ be events and $S_j=\sum_{i_1<\dots<i_j}P(A_{i_1}\cap\dots\cap A_{i_j})$.
\begin{enumerate}[label=(\alph*)]
\item Show that $P(A_1\cup\dots\cup A_m)\le\sum_{j=1}^t(-1)^{j-1}S_j$ for $t$ odd, and $\ge$ for $t$ even.
\item Deduce that the probability $D_n/n!$ that a random permutation of $[n]$ has no fixed point differs from $\sum_{j=0}^t\frac{(-1)^j}{j!}$ by at
most $\frac1{(t+1)!}$, for $0\le t\le n$.
\end{enumerate}
\end{exercise}

\begin{exercise}[$\star\star\star$]\label{pb:buchstab}
\emph{Rough numbers and Buchstab's function.}
\begin{enumerate}[label=(\alph*)]
\item Show that $\Phi(x,y)=1+\pi(x)-\pi(y)$ for $\sqrt x\le y\le x$, and that for $x^{1/3}\le y<\sqrt x$,
\[
\Phi(x,y)=1+\pi(x)-\pi(y)+\#\{(p,q):y<p\le q,\ pq\le x\}.
\]
\item Let $1<u\le3$. Show that $\frac{\log y}x\Phi(x,y)\to\omega(u)$ when $y=x^{1/u}$, where $\omega(u)=\frac1u$ for $1<u\le2$ and $\omega(u)=\frac{1+\log(u-1)}u$ for $2\le u\le3$.
\item Check that $(u\omega(u))'=\omega(u-1)$ for $2<u<3$ (Buchstab's equation, which defines $\omega$ for all $u>1$). Show that $\omega(2)<e^{-\gamma}<\omega(3)$.
\item Independent sieving by the primes up to $y$ would predict $\Phi(x,y)\approx x\prod_{p\le y}(1-\frac1p)\sim e^{-\gamma}\frac x{\log y}$, that is $\omega\equiv e^{-\gamma}$. Where on $]1,3]$ is $\omega$ smallest, and by what factor
does the prediction overestimate there? Relate this to Legendre's failure (Exercise~8.2).
\end{enumerate}
\end{exercise}

\begin{figure}[ht]
\centering
\fig{ch13-buchstab}
\caption{Buchstab's function $\omega(u)$ (charcoal), computed from $\omega(u)=\frac1u$ on $[1,2]$ and Buchstab's equation, with the proportions $\frac{\log y}x\Phi(x,y)$ for
$x=10^8$ and $y=x^{1/u}$, $u=1.5,2,2.5,3,4$ (ochre), and the level $e^{-\gamma}$ predicted by independent sieving (rust, dashed). The function crosses $e^{-\gamma}$
between $2$ and $3$, again near $3.2$ and $4$, and is within $10^{-4}$ of it from $u=4$ on; the data are off by a term of order $\frac1{\log y}$ (Problem~\ref{pb:buchstab}).}
\label{fig:buchstab}
\end{figure}

\section{Hints}

\begin{hint}{ex:legendrepi}
$2^{\pi(z)}\le2^z$, and $\prod_{p\le z}(1-\frac1p)=O(\frac1{\log z})$.
\end{hint}
\begin{hint}{ex:legendrefail}
Subtract Legendre's identity from its main term and use $\Phi(x,\sqrt x)=\pi(x)-\pi(\sqrt x)+1$.
\end{hint}
\begin{hint}{ex:sophie}
Sieve $\{n(2n+1):n\le x\}$, with $\rho(2)=1$ and $\rho(p)=2$ for $p$ odd.
\end{hint}
\begin{hint}{ex:btap}
Sieve $\{a+qm:m\le\frac{x-a}q\}$ by the primes not dividing $q$. Show $\sum_{d\le y,\,(d,q)=1}\frac{\mu^2(d)}{\varphi(d)}\ge\frac{\varphi(q)}q\log y$ by multiplying by $\sum_{k\mid q}\frac{\mu^2(k)}{\varphi(k)}=\frac q{\varphi(q)}$.
\end{hint}
\begin{hint}{ex:goldbachup}
Sieve $\{n(N-n):n\le N\}$. Show that $\sum_{a\le Y,\,(a,N)=1}\frac1a\ge\frac{\varphi(N)}N\log Y$.
\end{hint}
\begin{hint}{ex:twinheuristic}
(a) Abel summation as in Corollary~\ref{cor:brun}. (b) Compute $\frac{1-2/p}{(1-1/p)^2}$ and $\frac{1-1/2}{(1-1/2)^2}$. (c) $\int_x^\infty\frac{2\,\mathrm d\pi_2(t)}t$.
\end{hint}
\begin{hint}{ex:lsprimes}
$\omega(p)=1$ for $p\le Q$, so $L=\sum_{q\le Q}\frac{\mu^2(q)}{\varphi(q)}$. Remove the primes $\le Q$ first.
\end{hint}
\begin{hint}{ex:nsq1}
$\rho(p)=1+\chi_{-4}(p)$ for odd $p$. For the lower bound of $G$, compare with $\sum_{m\le y}\frac{(1*\chi_{-4})(m)}m=L(1,\chi_{-4})\log y+O(1)$, which follows from Dirichlet's hyperbola method (Lemma~\ref{lem:hyperbola}).
\end{hint}

\begin{hint}{pb:bonferroni}
A point lying in exactly $r$ of the events contributes $\sum_{j\le t}(-1)^{j-1}\binom rj=1-(-1)^t\binom{r-1}t$.
\end{hint}

\begin{hint}{pb:buchstab}
(a) A $y$-rough $n\le x$ with $y\ge x^{1/3}$ has at most two prime factors. (b) Count the pairs as $\sum_{y<p\le\sqrt x}(\pi(\frac xp)-\pi(p)+1)$, and turn
$\sum_{x^a<p\le x^b}\frac1pF(\frac{\log p}{\log x})$ into $\int_a^bF(t)\frac{\mathrm dt}t$ with Mertens's second theorem.
\end{hint}

\section{Solutions}

\begin{solution}{ex:legendrepi}
\idea Sieve only by the primes up to a tiny $z$, so that the $2^{\pi(z)}$ error is harmless, and accept the weak main term $\frac x{\log z}$.
\writeup Every prime $p\le x$ is either $\le z$ or has no prime factor $\le z$, so $\pi(x)\le z+\Phi(x,z)\le z+x\prod_{p\le z}(1-\frac1p)+2^z$. With $z=\log x$, $2^z=x^{\log2}$ and, by Theorem~\ref{thm:mertens3},
$\prod_{p\le z}(1-\frac1p)=O(\frac1{\log z})=O(\frac1{\log\log x})$. All three terms are $O(\frac x{\log\log x})$.
\end{solution}

\begin{solution}{ex:legendrefail}
\idea Legendre's identity is exact, so the error is whatever separates the true count from the main term, and both are known asymptotically.
\writeup Legendre's identity says that the sum in question equals $x\prod_{p\le\sqrt x}(1-\frac1p)-\Phi(x,\sqrt x)$. By Theorem~\ref{thm:mertens3}, the first term is $\sim\frac{e^{-\gamma}x}{\log\sqrt x}=\frac{2e^{-\gamma}x}{\log x}$. The survivors are $1$ and
the primes in $]\sqrt x,x]$, so $\Phi(x,\sqrt x)=\pi(x)-\pi(\sqrt x)+1\sim\frac x{\log x}$ by the prime number theorem. The difference is $\sim(2e^{-\gamma}-1)\frac x{\log x}$, that is $0.1229\frac x{\log x}$. Numerical evaluation shows the convergence
is slow: the ratio of main term to count is $0.998$ at $x=10^4$ and $1.034$ at $10^6$, on its way to $1.1229$.
\end{solution}

\begin{solution}{ex:sophie}
\idea Exactly the twin prime sieve, with $2n+1$ in place of $n+2$: only the number of classes removed per prime matters.
\writeup Let $\mathcal A=\{n(2n+1):n\le x\}$. For odd $p$, $p\mid n(2n+1)$ for $n\equiv0$ and $n\equiv-\frac12\pmod p$, two classes; for $p=2$, only $n\equiv0$. So $\rho$ is the same as in Theorem~\ref{thm:twins}, and so are $g$,
$h$, the lower bound for $G(y)$ and the remainder. A Sophie Germain prime $p>z$ gives an element $p(2p+1)$ coprime to $P(z)$. The proof of Theorem~\ref{thm:twins} applies word for word.
\end{solution}

\begin{solution}{ex:btap}
\idea The progression is an interval in disguise, of length $\frac xq$, and the primes dividing $q$ are not to be sieved. Their absence costs exactly the factor $\frac q{\varphi(q)}$ in $G$.
\writeup Let $\mathcal A=\{a+qm:0\le m\le\frac{x-a}q\}$, $X=\frac xq$, and $\mathcal P$ the primes not dividing $q$. For $d$ squarefree and coprime to $q$, $d\mid a+qm$ is one class of $m$ modulo $d$, so
$|\mathcal A_d|=\frac Xd+r_d$, $|r_d|\le1$: $g(d)=\frac1d$, $h(d)=\frac1{\varphi(d)}$. Every squarefree $d\le y$ factors uniquely as $kd'$ with $k\mid q$ and $(d',q)=1$, and $\frac1{\varphi(d)}=\frac1{\varphi(k)\varphi(d')}$. Hence
\[
\begin{aligned}\frac q{\varphi(q)}\sum_{\substack{d'\le y\\(d',q)=1}}\frac{\mu^2(d')}{\varphi(d')}&=\sum_{k\mid q}\frac{\mu^2(k)}{\varphi(k)}\sum_{\substack{d'\le y\\(d',q)=1}}\frac{\mu^2(d')}{\varphi(d')}\eqbr\ge\sum_{d\le y}\frac{\mu^2(d)}{\varphi(d)}\ge\log y\end{aligned}
\]
by Lemma~\ref{lem:Glog}, using $\sum_{k\mid q}\frac{\mu^2(k)}{\varphi(k)}=\prod_{p\mid q}(1+\frac1{p-1})=\frac q{\varphi(q)}$. Take $y=z=\frac{\sqrt X}{\log X}$. Then
\[
\begin{aligned}\pi(x;q,a)&\le z+\frac{X}{G(y)}+y^2\\&\le z+\frac q{\varphi(q)}\frac{x/q}{\frac12\log X-\log\log X}\eqbr+\frac X{\log^2X},\end{aligned}
\]
which is the claim since $\frac{X}{\log^2X}$ and $z$ are $O\bigl(\frac{x}{\varphi(q)\log^2(x/q)}\bigr)$, using $\varphi(q)\le q$.
\end{solution}

\begin{solution}{ex:goldbachup}
\idea Sieve $n$ and $N-n$ at once. A prime dividing $N$ divides $n$ and $N-n$ together, so it removes one class instead of two; that is where $\frac N{\varphi(N)}$ comes from.
\writeup Let $\mathcal A=\{n(N-n):1\le n\le N\}$. For $p\nmid N$, $p\mid n(N-n)$ for two classes of $n$; for $p\mid N$, one class. So $g(p)=\frac2p$ or $\frac1p$, and $h(p)=\frac2{p-2}$ or $\frac1{p-1}$. Restricting $G$ to odd $d$ coprime to $N$ and arguing as
in Theorem~\ref{thm:twins}, $G(y)\ge\bigl(\sum_{a\le\sqrt y,\,(a,2N)=1}\frac1a\bigr)^2$. Every $n\le Y$ is uniquely $kb$ with $k$ composed of primes dividing $2N$ and $(b,2N)=1$, $b\le Y$, so
$\prod_{p\mid2N}(1-\frac1p)^{-1}\sum_{b\le Y,(b,2N)=1}\frac1b\ge\sum_{n\le Y}\frac1n\ge\log Y$. As $N$ is even, $\frac{\varphi(2N)}{2N}=\frac{\varphi(N)}N$, so $G(y)\ge\bigl(\frac{\varphi(N)}N\bigr)^2\frac{\log^2y}4$. The remainder is bounded as for twins. With
$y=z=\frac{\sqrt N}{\log^3N}$, $r(N)\le z+\frac N{G(y)}+\frac N{\log^4N}$, which is the claim. (Since $\frac N{\varphi(N)}=O(\log\log N)$, this is $O\bigl(\frac{N(\log\log N)^2}{\log^2N}\bigr)$; Hardy and Littlewood predicted $r(N)\sim2C_2\prod_{p\mid N,\,p>2}\frac{p-1}{p-2}\cdot\frac N{\log^2N}$.)
\end{solution}

\begin{solution}{ex:twinheuristic}
\idea A sum over a thin set converges as fast as the set thins: with density $\frac1{\log^2t}$, the tail of $\sum\frac1t$ is $\int_x^\infty\frac{\mathrm dt}{t\log^2t}=\frac1{\log x}$.
\writeup (a) By Abel summation, $\sum_{p>x}\frac1p\le\int_x^\infty\frac{\pi_2(t)}{t^2}\dd t=O\bigl(\int_x^\infty\frac{\mathrm dt}{t\log^2t}\bigr)=O(\frac1{\log x})$ (the boundary term $-\frac{\pi_2(x)}x$ is negative), and the same for $\frac1{p+2}<\frac1p$.

(b) For an odd prime $p$, independence would give $\Pr(p\nmid n)\Pr(p\nmid n+2)=(1-\frac1p)^2$, while the true proportion of $n$ with $p\nmid n(n+2)$ is $1-\frac2p$. For $p=2$: $(1-\frac12)^2=\frac14$ against the true $\frac12$. The
correction factor is $\frac{1/2}{1/4}\prod_{p>2}\frac{1-2/p}{(1-1/p)^2}=2\prod_{p>2}\frac{p(p-2)}{(p-1)^2}=2\prod_{p>2}\Bigl(1-\frac1{(p-1)^2}\Bigr)=2C_2$. The product converges absolutely. This is the same
correction, prime by prime, that Chapter~1 found missing in the naive sieve heuristic; here it is applied only to the ratio of two sieve problems, where the troublesome factor $2e^{-\gamma}$ cancels.

(c) Abel summation with $\pi_2(t)\sim\frac{2C_2t}{\log^2t}$ gives $\sum_{p>x}\frac1p\sim\int_x^\infty\frac{2C_2}{t\log^2t}\dd t=\frac{2C_2}{\log x}$, and the same for $\frac1{p+2}$: total $\frac{4C_2}{\log x}$. At $x=10^7$, $\frac{4C_2}{\log x}=0.16383$, and
$1.7383570+0.1638312=1.90219$, against $B_2=1.9021605\ldots$ The same estimate from $10^5$ gives $1.90216$: the convergence of Brun's series itself is hopelessly slow, but the extrapolation is excellent.
\end{solution}

\begin{solution}{ex:lsprimes}
\idea Primes avoid one class, $0$, modulo every small prime, which is exactly what the large sieve needs, and then $L$ is the same sum $G$ as in Selberg's sieve.
\writeup Let $Q=\frac{\sqrt N}{\log N}$ and $\mathcal A$ the primes in $]M,M+N]$ that are $>Q$. They avoid the class $0$ modulo each $p\le Q$, so $\omega(p)=1$, $\frac{\omega(p)}{p-\omega(p)}=\frac1{p-1}$, $L=\sum_{q\le Q}\frac{\mu^2(q)}{\varphi(q)}\ge\log Q$. Theorem~\ref{thm:arithls} gives
$|\mathcal A|\le\frac{\pi N+Q^2}{\log Q}=\frac{\pi N(1+o(1))}{\frac12\log N}$, and the primes $\le Q$ add $O(\sqrt N)$. The loss is entirely in the constant of Theorem~\ref{thm:largesieve}: with the sharp $N-1+Q^2$ the large sieve gives $\frac{2N}{\log N}$,
exactly like Selberg's sieve. That the two methods agree is no accident: Selberg's sieve and the large sieve are dual to each other.
\end{solution}

\begin{solution}{ex:nsq1}
\idea The polynomial $n^2+1$ has $\rho(p)$ roots modulo $p$, and $\rho(p)$ is $2$ or $0$ according to $p\bmod4$; on average it is $1$, so we sieve ``one class per prime'', and $G(y)$ grows like $\log y$, with the constant $L(1,\chi_{-4})=\frac\pi4$.
\writeup Let $\mathcal A=\{n^2+1:n\le x\}$ and $\mathcal P=\{2\}\cup\{p\equiv1\ (4)\}$; primes $p\equiv3\pmod4$ never divide $n^2+1$ (Volume~3), so sieving by them is pointless. For $p\in\mathcal P$ let $\rho(p)$ be the number of roots of $n^2+1$ modulo
$p$: $\rho(2)=1$, $\rho(p)=2$. Then $|\mathcal A_d|=\frac{\rho(d)}dx+r_d$, $|r_d|\le\rho(d)\le\tau(d)$, and the remainder is handled as in Theorem~\ref{thm:twins}. As there, for squarefree odd $d$ composed of primes $\equiv1\pmod4$,
$h(d)\ge\sum_{\operatorname{rad}(m)=d}\frac{\tau(m)}m$, so $G(y)\ge\sum'_{m\le y}\frac{\tau(m)}m$, the sum over $m$ composed of primes $\equiv1\pmod4$.

Now compare with $f=1*\chi_{-4}$, which is multiplicative with $f(p^k)=k+1$ for $p\equiv1$, $f(2^k)=1$, and $f(p^k)=1$ or $0$ ($k$ even or odd) for $p\equiv3\pmod4$. So every $m\le y$ with $f(m)\ne0$ is $m'k$ with $m'$ composed of primes
$\equiv1\pmod4$, $f(m)=\tau(m')$, and $k=2^jl^2$ with $l$ composed of primes $\equiv3\pmod4$. Hence
\[
\begin{aligned}\sum_{m\le y}\frac{f(m)}m&\le\sum'_{m'\le y}\frac{\tau(m')}{m'}\sum_{j\ge0}\frac1{2^j}\sum_{l\ge1}\frac1{l^2}\eqbr=2\zeta(2)\sum'_{m'\le y}\frac{\tau(m')}{m'}.\end{aligned}
\]
By the hyperbola method, $\sum_{m\le y}\frac{f(m)}m=\sum_{ab\le y}\frac{\chi_{-4}(b)}{ab}=L(1,\chi_{-4})\log y+O(1)$, with $L(1,\chi_{-4})=\frac\pi4$ (Chapter~3). So $G(y)\ge\frac{\pi/4}{2\zeta(2)}\log y(1+o(1))=\frac3{4\pi}\log y\,(1+o(1))$.
Take $y=z=\frac{\sqrt x}{\log^2x}$: the remainder is at most $\bigl(y(1+\log y)\bigr)^2=O(\frac x{\log^2x})$, and
$\#\{n\le x:n^2+1\text{ prime}\}\le z+\frac x{G(y)}+O(\frac x{\log^2x})=O(\frac x{\log x})$. (The conjectured count is $\sim C\frac x{\log x}$ with $C=\frac12\prod_{p>2}\bigl(1-\frac{\chi_{-4}(p)}{p-1}\bigr)=0.6864\ldots$, the
Hardy--Littlewood constant for $n^2+1$; the sieve gets the right order of magnitude, and the parity problem forbids more.)
\end{solution}

\begin{solution}{pb:bonferroni}
\idea Look at one outcome at a time. If it lies in $r$ of the events, the truncated alternating sum counts it $1-(-1)^t\binom{r-1}t$ times,
which is above or below $1$ depending on the parity of $t$.

\writeup (a) For an outcome $\omega$ in exactly $r\ge1$ events, the indicator of the union is $1$, and $\omega$ contributes $\binom rj$ to $S_j$ (it lies in
$\binom rj$ of the intersections of $j$ events). By induction on $t$ from Pascal's rule,
$\sum_{j=1}^t(-1)^{j-1}\binom rj=1-(-1)^t\binom{r-1}t$. This is $\ge1$ for $t$ odd and $\le1$ for $t$ even, and for $r=0$ every term is $0$.
Taking expectations of these pointwise inequalities gives (a).

(b) Let $A_i$ be the event that $i$ is a fixed point. An intersection of $j$ of them has probability $\frac{(n-j)!}{n!}$, so $S_j=\binom nj\frac{(n-j)!}{n!}=\frac1{j!}$.
Then $\frac{D_n}{n!}=1-P(\bigcup A_i)$ lies between $\sum_{j=0}^t\frac{(-1)^j}{j!}$ and $\sum_{j=0}^{t+1}\frac{(-1)^j}{j!}$, which differ by $\frac1{(t+1)!}$ (for $t<n$;
for $t=n$ it is exact inclusion--exclusion).

\tools Volume~1, Chapter~2 (inclusion--exclusion, derangements).
\leads Brun's sieve (1915) is a clever truncation of inclusion--exclusion over primes, using Bonferroni at a well-chosen $t$; it gave
the first upper bound for twin primes (Volume~6, Chapter~8).
\end{solution}

\begin{solution}{pb:buchstab}
\idea A rough number below $x$ with no prime factor below $x^{1/3}$ is $1$, a prime, or a product of two primes; counting the products with the prime number theorem gives
$\int\frac{\mathrm dt}{t(1-t)}$. Independent sieving would give $e^{-\gamma}$ everywhere, and the actual answer is far from it when $y$ is a large power of $x$.

\writeup (a) A $y$-rough $n\le x$ with $n>1$ and $y\ge x^{1/3}$ has at most two prime factors, since three would multiply to more than $y^3\ge x$. If $y\ge\sqrt x$ it cannot have two
either, so it is a prime in $]y,x]$. If $x^{1/3}\le y<\sqrt x$, it is a prime in $]y,x]$ or a product $pq$ with $y<p\le q$, $pq\le x$.

(b) For $1<u\le2$, $\Phi(x,y)=1+\pi(x)-\pi(y)$, which is $\sim\frac x{\log x}$ since $\pi(x^{1/u})=o(\pi(x))$. Multiplying by $\frac{\log y}x=\frac{\log x}{ux}$ gives $\frac1u$.
For $2<u\le3$, the pairs number
\[
\sum_{y<p\le\sqrt x}\Bigl(\pi\Bigl(\frac xp\Bigr)-\pi(p)+1\Bigr),
\]
which is $\sum_{y<p\le\sqrt x}\pi(\frac xp)+O(\pi(\sqrt x)^2)$.
Since $\frac xp\ge\sqrt x$, $\pi(\frac xp)=(1+o(1))\frac{x}{p\log(x/p)}$ uniformly, and the sum is
\[
(1+o(1))\frac x{\log x}\sum_{y<p\le\sqrt x}\frac1pF\Bigl(\frac{\log p}{\log x}\Bigr)\qquad\text{with }F(t)=\frac1{1-t}.
\] By Mertens's
second theorem (Theorem~1.5), $\sum_{x^a<p\le x^b}\frac1p\to\log\frac ba$ for fixed $0<a<b$; cutting $[\frac1u,\frac12]$ into small intervals and using the continuity of $F$, Riemann sums give
\[
\sum_{x^{1/u}<p\le x^{1/2}}\frac1pF\Bigl(\frac{\log p}{\log x}\Bigr)\to\int_{1/u}^{1/2}\frac{\mathrm dt}{t(1-t)}=\Bigl[\log\frac t{1-t}\Bigr]_{1/u}^{1/2}=\log(u-1).
\]
So $\Phi(x,y)=(1+\log(u-1)+o(1))\frac x{\log x}$ and $\frac{\log y}x\Phi(x,y)\to\frac{1+\log(u-1)}u$. (At $u=2$ both formulas give $\frac12$.)

(c) For $2<u<3$, $(u\omega(u))'=\frac1{u-1}=\omega(u-1)$, as $u-1\in\left]1,2\right[$. Numerically $\omega(2)=0.5<e^{-\gamma}=0.56146<\omega(3)=\frac{1+\log2}3=0.56438$.

(d) $\omega$ decreases from $1$ to $\frac12$ on $]1,2]$; on $[2,3]$, $\omega(u)\ge\frac12$ because $2+2\log(u-1)-u$ vanishes at $u=2$ and has derivative $\frac2{u-1}-1\ge0$. So the minimum on $]1,3]$ is
$\omega(2)=\frac12$, where independent sieving overestimates by the factor $2e^{-\gamma}=1.1229$. This is Legendre's failure: sieving the integers up to $x$ by the primes up to $\sqrt x$
leaves the primes of $]\sqrt x,x]$, about $\frac x{\log x}$ of them, not $x\prod_{p\le\sqrt x}(1-\frac1p)\sim2e^{-\gamma}\frac x{\log x}$ (Exercise~8.2). As $u\to1$, on the other hand, the prediction
underestimates by a factor tending to $e^\gamma$.

\tools{Theorems~1.5 and~5.5, Exercise~8.2.}
\leads Buchstab (1937) found the equation $(u\omega(u))'=\omega(u-1)$, and de Bruijn (1950) proved that $\omega(u)\to e^{-\gamma}$, the difference oscillating and decaying faster than any exponential
(Figure~\ref{fig:buchstab}). So independent sieving is right exactly when the primes used are a small power of the range: for $y=x^{1/u}$ with $u$ large. This is the ``sieve limit'' of Volume~6,
Chapter~8, seen from the exact side.
\end{solution}

\section*{Sources and further reading}

\begin{itemize}
\item A.-M.~Legendre, \emph{Essai sur la théorie des nombres}, 2nd ed., Paris, 1808.
\item V.~Brun, \emph{La série $\frac15+\frac17+\frac1{11}+\frac1{13}+\cdots$ où les dénominateurs sont nombres premiers jumeaux est convergente ou finie}, Bulletin des Sciences Mathématiques 43 (1919), 100--104 and 124--128.
\item A.~Selberg, \emph{On an elementary method in the theory of primes}, Det Kongelige Norske Videnskabers Selskabs Forhandlinger 19 (1947), 64--67.
\item Yu.~V.~Linnik, \emph{The large sieve}, Doklady Akademii Nauk SSSR 30 (1941), 292--294.
\item P.~X.~Gallagher, \emph{The large sieve}, Mathematika 14 (1967), 14--20.
\item H.~L.~Montgomery, \emph{A note on the large sieve}, Journal of the London Mathematical Society 43 (1968), 93--98.
\item H.~L.~Montgomery and R.~C.~Vaughan, \emph{The large sieve}, Mathematika 20 (1973), 119--134.
\item G.~H.~Hardy and J.~E.~Littlewood, \emph{Some problems of ``Partitio numerorum'' III}, Acta Mathematica 44 (1923), 1--70.
\item T.~R.~Nicely, \emph{Enumeration to $10^{14}$ of the twin primes and Brun's constant}, Virginia Journal of Science 46 (1995), 195--204.
\item J.-R.~Chen, \emph{On the representation of a larger even integer as the sum of a prime and the product of at most two primes}, Scientia Sinica 16 (1973), 157--176.
\item Y.~Zhang, \emph{Bounded gaps between primes}, Annals of Mathematics 179 (2014), 1121--1174; J.~Maynard, \emph{Small gaps between primes}, Annals of Mathematics 181 (2015), 383--413.
\item H.~Halberstam and H.-E.~Richert, \emph{Sieve Methods}, Academic Press, 1974. The classical reference.
\item E.~Bombieri and H.~Davenport, Small differences between prime numbers, \emph{Proc. Roy. Soc. London Ser. A} 293 (1966), 1--18 (the factor $4$ for twin primes). J.~Wu, Chen's double sieve, Goldbach's conjecture and the twin prime problem, \emph{Acta Arith.} 114 (2004), 215--273 (the factor $3.4$).
\item J.~Friedlander and H.~Iwaniec, \emph{Opera de Cribro}, AMS Colloquium Publications 57, 2010. The modern one, witty and complete.
\item H.~Davenport, \emph{Multiplicative Number Theory}, 3rd ed., Springer GTM~74, 2000, Chapters~23 and~27--28 (the large sieve, Bombieri--Vinogradov).
\end{itemize}
